% ============================================================================
% Section II: Related Work
% ============================================================================
\section{Related Work}
\label{sec:related_work}

\subsection{Fast-Charging Protocol Optimization}

Fast-charging design spans model-based optimal control, data-driven search, and
closed-loop experimentation. Attia \textit{et al.} combined early lifetime
prediction with BO in physical-cell experiments~\cite{ref11}. Jiang and
co-workers studied minimum-time charging, acquisition choices, protocol
dimensionality, and operational constraints~\cite{ref12,ref14}, while Wang and
Jiang used Chebyshev decomposition for constrained multiobjective charging
design~\cite{ref13}. These studies establish the relevance of sequential
surrogate optimization for expensive charging design and motivate the
three-objective formulation considered here.

Recent work accelerates this process through several sources of structure:
mesh-adaptive search~\cite{refDong}, online digital-twin
models~\cite{refDigitalTwin}, and parallel satisficing Thompson
sampling~\cite{refBLASTS}. Other studies use early-life prediction with noisy
expected hypervolume improvement~\cite{refZhu}, a semi-empirical aging-model
prior~\cite{refPrBO}, or section-wise electrochemical constraints with
physical-cell validation~\cite{refYoon}. These approaches derive guidance from
numerical models, related-task data, or evaluated cells. LLMBO-MO instead uses
textual context for screened initialization and regional preference while
retaining simulator evaluation as the sole source of objective data.

\subsection{Expensive Multiobjective Optimization}

ParEGO decomposes an expensive vector-valued problem into scalar subproblems
and applies a single-objective surrogate and acquisition rule~\cite{ref4}.
Constraint-aware BO can additionally model the probability of satisfying
expensive inequalities~\cite{refGardner}. Direct multiobjective alternatives
optimize expected hypervolume improvement, including differentiable qEHVI and
its noise-aware qNEHVI extension~\cite{refQEHVI,refQNEHVI}. Surrogate-assisted
evolutionary algorithms such as DISK and PIMD instead use Kriging predictions
with distribution- or infill-based selection~\cite{refDISK,refPIMD}; their
implementations are available in PlatEMO~\cite{refPlatEMO}. The present method
uses a ParEGO-style decomposition so that each iteration needs one scalar GP
and the fixed-budget comparison retains a common numerical backbone.

\subsection{LLM-Assisted Bayesian Optimization}

LLAMBO demonstrated that LLMs can support zero-shot initialization, surrogate
reasoning, and candidate sampling in data-sparse BO~\cite{ref10}. LGBO
introduced point and region preferences that alter a surrogate mean while
preserving its posterior covariance~\cite{ref16}. LABO takes a multi-fidelity
view, combining inexpensive LLM predictions with real-fidelity observations
through discrepancy modeling and a gating rule~\cite{refLABO}. LLMBO-MO is
closest to the preference-guided line but differs in the authority assigned to
language-model output: it neither inserts LLM predictions as objective values
nor lets the LLM choose the final evaluation. It adapts this
preference-guided pattern to a battery-specific, protocol-domain-screened
interface with initial portfolio construction, regional acquisition restarts,
and a bounded early acquisition-time mean adjustment.
